%%%%%%%%%%%%%%%%%%%%%%%%%%%%%%%%%%%%%%%%%%%%%%%%%%%%%%%%%%%%%%%%%
% GIÁO ÁN STEM: WAVEFORM - GIAI ĐIỆU TOÁN HỌC
% Môn: Toán 11 | Chủ đề: Hàm lượng giác & Sóng âm thanh
%%%%%%%%%%%%%%%%%%%%%%%%%%%%%%%%%%%%%%%%%%%%%%%%%%%%%%%%%%%%%%%%%
\documentclass[12pt,a4paper]{article}

\usepackage[utf8]{inputenc}
\usepackage[T1]{fontenc}
\usepackage[vietnamese]{babel}
\usepackage{graphicx}
\usepackage{amsmath,amssymb}
\usepackage{array, booktabs, tabularx}
\usepackage{tikz}
\usetikzlibrary{calc}
\usepackage{pgfplots}
\pgfplotsset{compat=1.18}
\usepackage{geometry}
\usepackage{fancyhdr}
\usepackage[svgnames]{xcolor}
\usepackage{tcolorbox}
\tcbuselibrary{skins, breakable}
\usepackage{hyperref}
\usepackage{enumitem}

\geometry{left=2cm,right=2cm,top=2.5cm,bottom=2cm}
\definecolor{wavepurple}{RGB}{139, 92, 246}    % Violet 500
\definecolor{wavedark}{RGB}{76, 29, 149}       % Violet 800
\definecolor{wavebg}{RGB}{245, 243, 255}       % Violet 50

\hypersetup{colorlinks=true, linkcolor=wavedark, urlcolor=wavepurple}

% ============= BOX DEFINITIONS =============
\newtcolorbox{sectionbox}[1]{
    enhanced, breakable, colback=white, colframe=wavepurple, boxrule=1pt,
    fonttitle=\bfseries\sffamily\Large, coltitle=white, colbacktitle=wavepurple,
    title={#1},
    shadow={2mm}{-1mm}{0mm}{black!20},
    code={\phantomsection\addcontentsline{toc}{section}{#1}}
}

\newtcolorbox{activitybox}[1]{
    enhanced, breakable, colback=white, colframe=wavedark, boxrule=1pt, arc=3mm,
    title={\sffamily\bfseries\color{white} [HOẠT ĐỘNG] #1},
    colbacktitle=wavedark
}

\newtcolorbox{mathbox}[1]{
    enhanced, breakable, colback=wavebg, colframe=wavedark, boxrule=1pt, arc=3mm,
    title={\sffamily\bfseries\color{white} [TOÁN] #1},
    colbacktitle=wavedark
}

% ============= HEADERS =============
\pagestyle{fancy}
\fancyhf{}
\fancyhead[L]{\color{wavedark}\textbf{DỰ ÁN STEM: WAVEFORM}}
\fancyhead[R]{\color{wavedark}\textbf{TOÁN HỌC \& ÂM NHẠC}}
\fancyfoot[C]{\thepage}

\begin{document}

% ========== TITLE PAGE ==========
\begin{titlepage}
    \begin{tikzpicture}[remember picture, overlay]
        \fill[wavedark] (current page.north west) rectangle (current page.south east);
        \node[anchor=center, color=wavepurple, font=\fontsize{60}{70}\sffamily\bfseries] at ($(current page.center)+(0,1)$) {WAVE};
        \node[anchor=center, color=white, font=\fontsize{60}{70}\sffamily\bfseries] at ($(current page.center)+(0,-1)$) {FORM};
        \node[anchor=center, color=white, font=\Large\sffamily] at ($(current page.center)+(0,-3)$) {Khi Âm nhạc Gặp Toán học};
        
        % Draw sine wave decoration
        \draw[wavepurple, line width=2pt, domain=0:360, samples=100] 
            plot ($(current page.south west)+(0.01*\x, 3 + 0.5*sin(\x))$);
    \end{tikzpicture}
    
    \vspace{12cm}
    \begin{center}
        \color{black}
        \begin{tcolorbox}[width=0.8\textwidth, colback=white, colframe=wavepurple]
            \centering \large
            \begin{tabular}{rl}
                \textbf{Môn học:} & Toán (Hàm lượng giác) \\
                \textbf{Chủ đề:} & Hàm Sin, Chu kỳ, Tần số \\
                \textbf{Ứng dụng:} & Âm thanh, Sóng, Nhạc cụ \\
                \textbf{Thời lượng:} & 2 tiết (90 phút)
            \end{tabular}
        \end{tcolorbox}
    \end{center}
\end{titlepage}

\tableofcontents
\newpage

% ========== CONTENT ==========

\begin{sectionbox}{I. MỤC TIÊU BÀI HỌC}
\phantomsection\subsection*{1. Kiến thức}\addcontentsline{toc}{subsection}{1. Kiến thức}
\begin{itemize}
    \item Nắm vững dạng tổng quát của hàm lượng giác: $y = A\sin(\omega x + \varphi)$.
    \item Hiểu ý nghĩa vật lý của các tham số: Biên độ ($A$), Tần số góc ($\omega$), Pha ban đầu ($\varphi$).
    \item Biết mối liên hệ: Chu kỳ $T = \frac{2\pi}{\omega}$, Tần số $f = \frac{1}{T}$.
\end{itemize}

\phantomsection\subsection*{2. Kỹ năng}\addcontentsline{toc}{subsection}{2. Kỹ năng}
\begin{itemize}
    \item Vẽ và phân tích đồ thị hàm sin với các tham số khác nhau.
    \item Sử dụng công cụ mô phỏng để "nghe" và "nhìn" sóng âm.
    \item Liên hệ tần số (Hz) với cao độ nốt nhạc.
\end{itemize}

\phantomsection\subsection*{3. Phẩm chất}\addcontentsline{toc}{subsection}{3. Phẩm chất}
\begin{itemize}
    \item Yêu thích vẻ đẹp của Toán học trong nghệ thuật.
    \item Tư duy liên môn, sáng tạo.
\end{itemize}
\end{sectionbox}

\begin{sectionbox}{II. CHUẨN BỊ}
\begin{itemize}
    \item \textbf{Giáo viên:} Máy chiếu, Loa, Web App \texttt{waveform.html}.
    \item \textbf{Học sinh:} Điện thoại/Laptop (có tai nghe), Phiếu học tập.
\end{itemize}
\end{sectionbox}

\begin{sectionbox}{III. TIẾN TRÌNH DẠY HỌC (MÔ HÌNH 5E)}

\phantomsection\subsection*{1. ENGAGE (Gắn kết - 10p): Âm thanh là gì?}\addcontentsline{toc}{subsection}{1. ENGAGE (Gắn kết - 10p)}
\textbf{Hoạt động:} GV phát 2 âm thanh (nốt Đồ và nốt La) qua loa.

\textbf{Câu hỏi:}
\begin{itemize}
    \item "Tại sao 2 nốt này nghe khác nhau?"
    \item "Điều gì quyết định một nốt nhạc cao hay thấp, to hay nhỏ?"
\end{itemize}

\textbf{Dẫn dắt:} "Âm thanh thực chất là Sóng. Và mọi sóng đều có thể mô tả bằng hàm Sin!"

\phantomsection\subsection*{2. EXPLORE (Khám phá - 20p): Nhìn thấy Âm thanh}\addcontentsline{toc}{subsection}{2. EXPLORE (Khám phá - 20p)}
\begin{activitybox}{Thực nghiệm trên WaveForm App}
Truy cập: \href{https://stem-1h8.pages.dev/waveform.html}{\texttt{stem-1h8.pages.dev/waveform.html}}

HS thực hiện:
\begin{enumerate}
    \item Bấm các phím đàn piano ảo, quan sát sóng sin hiển thị.
    \item So sánh sóng của nốt Đồ (C4 = 262 Hz) và nốt La (A4 = 440 Hz):
    \begin{itemize}
        \item Sóng nào "dày đặc" hơn (tần số cao hơn)?
        \item Sóng nào "thưa" hơn (tần số thấp hơn)?
    \end{itemize}
    \item Điều chỉnh thanh trượt Biên độ ($A$), quan sát sự thay đổi độ to của âm.
\end{enumerate}
\end{activitybox}

\phantomsection\subsection*{3. EXPLAIN (Giải thích - 25p): Toán học của Âm thanh}\addcontentsline{toc}{subsection}{3. EXPLAIN (Giải thích - 25p)}
\begin{mathbox}{Hàm Sóng Sin}
Một sóng âm thanh đơn giản (Pure Tone) được mô tả bởi:
$$ y(t) = A \sin(2\pi f t) $$

Trong đó:
\begin{itemize}
    \item $A$: \textbf{Biên độ} (Amplitude) - Độ to của âm thanh.
    \item $f$: \textbf{Tần số} (Frequency, Hz) - Số dao động trong 1 giây. Quyết định cao độ.
    \item $t$: Thời gian (giây).
    \item $T = \frac{1}{f}$: \textbf{Chu kỳ} - Thời gian hoàn thành 1 dao động.
\end{itemize}

\textbf{Bảng Tần số các Nốt nhạc (Octave 4):}
\begin{center}
\begin{tabular}{|c|c|c|c|c|c|c|c|}
\hline
Nốt & C4 & D4 & E4 & F4 & G4 & A4 & B4 \\
\hline
Tần số (Hz) & 262 & 294 & 330 & 349 & 392 & 440 & 494 \\
\hline
\end{tabular}
\end{center}
\end{mathbox}

\begin{center}
\begin{tikzpicture}
\begin{axis}[
    width=12cm, height=5cm,
    domain=0:0.02, samples=500,
    xlabel={Thời gian (s)}, ylabel={Biên độ},
    legend pos=outer north east,
    grid=major
]
\addplot[blue, thick] {sin(deg(2*pi*262*x))};
\addlegendentry{C4 (262 Hz)}
\addplot[red, thick] {sin(deg(2*pi*440*x))};
\addlegendentry{A4 (440 Hz)}
\end{axis}
\end{tikzpicture}
\end{center}

\phantomsection\subsection*{4. ELABORATE (Áp dụng - 25p): Sáng tác Giai điệu}\addcontentsline{toc}{subsection}{4. ELABORATE (Áp dụng - 25p)}
\begin{activitybox}{Thử thách Nhạc sĩ Toán học}
Sử dụng App để:
\begin{enumerate}
    \item Tìm tần số của nốt Đồ5 (C5). Gợi ý: Gấp đôi C4.
    \item Vẽ đồ thị $y = 2\sin(2\pi \cdot 330 \cdot t)$ (Nốt Mi với biên độ 2).
    \item Thử cộng 2 sóng sin để tạo hợp âm (Chord). Quan sát hiện tượng "Beat".
\end{enumerate}
\end{activitybox}

\phantomsection\subsection*{5. EVALUATE (Đánh giá - 10p)}\addcontentsline{toc}{subsection}{5. EVALUATE (Đánh giá - 10p)}
HS hoàn thành Phiếu học tập:
\begin{enumerate}
    \item Xác định tần số và chu kỳ của một nốt nhạc cho trước.
    \item Giải thích tại sao quãng tám (Octave) có tần số gấp đôi.
\end{enumerate}
\end{sectionbox}

\newpage
\begin{sectionbox}{IV. PHỤ LỤC}
\phantomsection\subsection*{Phụ lục 1: Ứng dụng thực tế}\addcontentsline{toc}{subsection}{Phụ lục 1: Ứng dụng thực tế}
\begin{itemize}
    \item Equalizer (EQ) trong phần mềm nhạc.
    \item Chẩn đoán y tế bằng sóng siêu âm.
    \item Công nghệ nhận dạng giọng nói (AI Voice).
\end{itemize}

\phantomsection\subsection*{Phụ lục 2: Mở rộng}\addcontentsline{toc}{subsection}{Phụ lục 2: Mở rộng}
\begin{itemize}
    \item Chuỗi Fourier: Mọi sóng phức tạp đều là tổng của các sóng sin đơn giản.
    \item Synthesizer: Tổng hợp âm thanh điện tử bằng Toán học.
\end{itemize}
\end{sectionbox}

\end{document}
